\section{Problemstellung und Zweck}
\label{sec:01-01_problem-and-context}

Digitale Wissens- und Kollaborationsplattformen wie Intranets und Wikis sind in Organisationen zentrale Infrastrukturen zur Bündelung von Informationen, zur Prozessdokumentation und zur bereichsübergreifenden Zusammenarbeit. Empirische Studien zeigen jedoch, dass diese Systeme ihr Potenzial häufig nicht ausschöpfen, wenn \Gls{ia} und \Gls{ux} nicht systematisch an Aufgaben, Nutzungskontext und \gls{mentalemodelle} der Mitarbeitenden ausgerichtet werden \cite{402:Rosenfeld2015,500:ISO9241-210}. Die Folgen sind erhöhte Suchzeiten, ineffiziente \Glspl{workaround} sowie eine geringe Akzeptanz der Plattformen als \enquote{\Gls{singlesourceoftruth}} \cite{300:FraunhoferIESE2018,301:Infragistics2015}.

\Gls{ux} wird in Normen und Fachliteratur als Gesamtheit der Wahrnehmungen und Reaktionen beschrieben, die aus der Nutzung eines Systems entstehen, einschließlich vorangehender Erwartungen sowie unmittelbarer und nachgelagerter Erlebnisse \cite{500:ISO9241-210}. Norman verdeutlicht, dass wahrgenommene Bedienfehler in komplexen Systemen meist auf unklare \Glspl{affordanz}, fehlendes \Gls{feedback} und inkonsistente Abbildungen zwischen Handlung und Systemzustand zurückzuführen sind und damit primär gestalterische Ursachen haben \cite{400:Norman2013}. In internen Geschäftsanwendungen erhöhen solche Defizite Fehlerwahrscheinlichkeit, Frustration und \gls{cognitiveload} und wirken sich unmittelbar auf Effizienz, Qualität und Zufriedenheit der Mitarbeitenden aus \cite{300:FraunhoferIESE2018,400:Norman2013}.

\Gls{ia} wird nach Rosenfeld, Morville und Arango als Kunst und Wissenschaft der Strukturierung, Benennung und Verknüpfung von Informationen verstanden, um deren Auffindbarkeit und Verständlichkeit zu unterstützen. Sie umfasst neben Navigationsstrukturen auch \Glspl{klassifikationssysteme}, \Glspl{metadaten}, Suchstrategien und kontextuelle Verlinkungen, die auf das Informationssuchverhalten der Zielgruppen abgestimmt sein müssen. In Unternehmensintranets führt eine historisch gewachsene, am Organigramm orientierte \Gls{ia} häufig zu \Glspl{informationssilos}, Redundanzen und unklaren Bezeichnungen, wodurch Mitarbeitende Informationen nur mit hohem Aufwand finden oder eigene Wissensstrukturen aufbauen \cite{402:Rosenfeld2015}.

Garrett beschreibt \Gls{ux} als mehrschichtiges Modell, in dem Strategie, Umfang, Struktur (\Gls{ia}), \Gls{interaktionsdesign} und Oberfläche konsistent aufeinander abgestimmt sein müssen. \Gls{ia} fungiert dabei als zentrales Bindeglied, das funktionale Anforderungen und Inhalte in eine für Nutzende nachvollziehbare Ordnung überführt \cite{401:Garrett2011}. Fehlende Abstimmung zwischen diesen Ebenen führt zu Brüchen zwischen organisationalen Zielen, Nutzungsanforderungen und tatsächlicher Interaktion, was sich insbesondere in komplexen Wissensplattformen negativ auswirkt \cite{401:Garrett2011,402:Rosenfeld2015}.

Die Norm \textit{ISO~9241-210} bündelt diese Perspektiven im Rahmen des \Gls{humancentereddesign} und fordert eine systematische Analyse von Nutzenden, Aufgaben und Nutzungskontexten, die kontinuierliche Einbeziehung der Nutzenden sowie iterative Evaluation von Gestaltungslösungen \cite{500:ISO9241-210}. Für interne Systeme bedeutet dies, dass \Gls{ia} und \Gls{ux} als fortlaufender, in Prozesse und \Gls{governance}-Strukturen integrierter Gestaltungsansatz zu verstehen sind, der organisatorische und technische Rahmenbedingungen berücksichtigt \cite{500:ISO9241-210,300:FraunhoferIESE2018}. Studien zeigen, dass \Gls{ux} in Geschäftsanwendungen direkten Einfluss auf Prozessqualität, Fehlerquote, Schulungsaufwand und Akzeptanz digitaler Lösungen hat und damit auch wirtschaftlich relevant ist \cite{300:FraunhoferIESE2018,301:Infragistics2015,100:WhartonUX2025}.

Im Zuge der \gls{digitaleTransformation} steigt die Komplexität interner Systemlandschaften durch verteilte, cloudbasierte Strukturen, \gls{hybridarbeit} und eine wachsende Anzahl digitaler Werkzeuge \cite{302:GSMA2024}. Untersuchungen zu \Gls{enterpriseux} zeigen, dass gezielte Investitionen in \Gls{ux} und informationsarchitektonische Klarheit Produktivität, Qualität und Nutzerzufriedenheit verbessern, während schlecht nutzbare Systeme \Gls{medienbruch}, ineffiziente Parallelstrukturen und Widerstände gegen Veränderungsprozesse fördern \cite{101:Khan2021,300:FraunhoferIESE2018,301:Infragistics2015}. Damit entwickelt sich \Gls{ux} im organisationalen Kontext von einem rein gestalterischen Aspekt zu einem strategischen Erfolgsfaktor für \Gls{wissensmanagement}, Kollaboration und organisationale Veränderungsfähigkeit \cite{100:WhartonUX2025,300:FraunhoferIESE2018}.

Zusammenfassend lässt sich festhalten, dass Organisationen zwar erheblich in digitale Wissens- und Kollaborationsplattformen investieren, deren \Gls{ia} und \Gls{ux} jedoch häufig nicht konsequent an Nutzerbedürfnissen, Arbeitsaufgaben und strategischen Transformationszielen ausrichten. Dies führt zu begrenzter Nutzung, Akzeptanzbarrieren und einer eingeschränkten Unterstützung von Wissensaustausch und Entscheidungsfindung \cite{301:Infragistics2015,402:Rosenfeld2015}. Die vorliegende Arbeit setzt an dieser Diskrepanz an und untersucht, wie systematisch gestaltete \Gls{ux}-Strategien und \Gls{ia} zur Wirksamkeit interner Plattformen in Organisationen beitragen können.

\newpage
